%%%%%%%%%%%%%%%%%%%%%%%%%%%%%%%%%%%%%%%%%%%%%%%%%%%%%%%%%%%%%%%%%%%%%%%%%%%
\chapter{Introduction}
%%%%%%%%%%%%%%%%%%%%%%%%%%%%%%%%%%%%%%%%%%%%%%%%%%%%%%%%%%%%%%%%%%%%%%%%%%%
\section{Project Context}

Studies in cognitive load in relation to neurodiversity and their influence in computer science and user experience mostly approach the topic in a qualitative manner, often using methodologies like literature review \cite{burtscher2024, hart2006, kosch2023, liu2024, skulmowski2022, spiel2022, stefanidi2022, wollberg2025, cognitive2025}, interviews \cite{hall2024, lowy2023}, and surveys \cite{conrad2023, robertson2024}. There exist some gaps in the methodologies that aim to quantify and represent cognitive load in HCI researches, especially with neurodivergent participants \cite{lecunff2024}. This study shall implement methodologies that supplement the lack of empirical approaches and potentially invaluable discoveries.

\section{Purpose and Description}

Due to the identified gaps in methodologies and target population of past researches, this study aims to empirically delineate the difference between the cognitive load in neurodivergent and neurotypical population. This study shall utilize the optimal tools and equipment for measuring cognitive load \cite{ahmad2023, appel2019, babu2018, baceviciute2020, hart2006, kosch2023, pretty2024} in these populations, anticipating the potential generation of new insights and knowledge. We shall also discover how the neurodivergent experience differs from the neurotypical experience, understanding the behaviors \cite{hall2024, lowy2023} and strengths \cite{stefanidi2022} of the former in relation to computer studies and systems design. Subsequently, this study will also highlight the lack of inclusivity \cite{lowy2023, wollberg2025} and accessibility \cite{elglaly2024, kuang2024, le2024} to applications and technologies that supposedly cater to neurodivergent populations.

\section{Objectives}

This study aims to:
\begin{itemize}
    \item Get a quantified representation or amount of cognitive load from neurodivergent students in comparison to neurotypical students while using the most common websites.
    \item Understand the experiences of neurodivergent students in relation to HCI and computing by assessing their cognitive load.
    \item Provide a theoretical framework and supporting study for projects that cater to the neurodivergent population.
\end{itemize} 

\section{Scope and Limitations}
This study only aims to look at the cognitive load of neurodivergent and neurotypical students of the same age and education level at the same educational institution. Cognitive stimulation is limited to the use of software to understand the user experience of these user groups. The research will focus on collecting physiological data to monitor the cognitive load of students. The study will not look into the performance of the students while doing various tasks.

\section{Significance}
Understanding the differences of neurotypical and neurodivergent cognitive load under software usage can have a significant impact on future software system designs. This can become a theoretical basis for designing universal and accessible software to allow neurodivergent individuals to benefit from modern software. This study can help fill the gaps of research around the neurodivergent experience \cite{lecunff2024}. The results can serve as a guideline for people involved in software development including the UI/UX Developers, Frontend Developers, and Software Testers.